\documentclass[11pt]{article}
\usepackage{elicitation}
\begin{document}

\packettitle{The Elicitation Round}{What to do, and why it is the last thing outstanding}

\begin{note}
This packet is everything needed to run the one open item on the project. Nothing in the
repository is blocked on engineering. This is the only task left that could change a conclusion
rather than qualify one.
\end{note}

\section{Why it matters}

Part Four of the book rests on a twelve-by-eight table that one person wrote down. Every
robustness check in the project varies the \textit{values} in that table and holds its pattern of
empty cells fixed, so no amount of further computation can test the pattern itself. The book
already says as much, and records that Part Four's coupling claims largely do not survive
structural randomisation. A second reader filling the same grid independently is the only test
available.

The result is genuinely two-sided. If two or three readers leave the same rows empty, a central
claim has independent support for the first time. If they do not, \textbf{Chapter 18 is wrong},
and the preregistered script says so in those words rather than softening it.

\section{The five documents}

\begin{center}
\begin{tabular}{@{}ll@{}}
\toprule
\textbf{Document} & \textbf{Who sees it} \\
\midrule
\texttt{0-start-here} & you \\
\texttt{1-respondent-form} & \textbf{the respondent} --- this is the document you send \\
\texttt{2-recruiting-note} & you, to copy into an email or message \\
\texttt{3-response-template} & send with the form, or let them return the form filled in \\
\texttt{4-collator-notes} & \textbf{you only} \\
\bottomrule
\end{tabular}
\end{center}

\begin{warning}
\textbf{Do not send \texttt{4-collator-notes} to a respondent.} It states how many rows the book
leaves empty and which they are, which is the answer the exercise is trying to elicit
independently. Sending it destroys the round.
\end{warning}

\section{Steps}

\begin{enumerate}[leftmargin=*,itemsep=6pt]
\item \textbf{Recruit two or three people.} Three is the useful minimum: with three you can
      report, per cell, whether zero, one, two or three marked it, and the empty-rows claim
      becomes a statement about agreement rather than about one person.

      The ideal respondent knows AA well enough to have opinions about what the Traditions do, and
      has \textit{not} read Part Four. Someone who has read it is not independent; that does not
      disqualify them, but it must be recorded, and their form reported separately.

\item \textbf{Send them \texttt{1-respondent-form}} and the note in \texttt{2-recruiting-note}.
      Twenty to forty minutes of work.

\item \textbf{Save each completed form} in the repository as
      \texttt{research/governance-elicitation-<initials>.md}. The script reads every file matching
      that pattern: a markdown table whose first column is the Tradition and whose remaining eight
      columns are the resources in the order given on the form. Blank and \texttt{0} both mean
      empty.

\item \textbf{Run the comparison:}

\begin{quote}\ttfamily\small
python3 model/elicitation\_compare.py -{}-self-test\\
python3 model/elicitation\_compare.py
\end{quote}

      Run the self-test before sending any forms out. It passed on 10 August 2026.

\item \textbf{Record who filled each one in}, what they were told, and whether they had seen the
      book, in \texttt{research/SOURCES.md}. Then propagate whatever the result says through the
      layers listed in \texttt{CLAUDE.md}, exactly as any other change to a public claim.
\end{enumerate}

\section{The instrument was checked before this packet was made}

\texttt{model/elicitation\_compare.py -{}-self-test} passed on 10 August 2026. It exercises two
synthetic respondents: one who matches the book's empty rows, and one who differs by a single row.
The second is correctly reported as contradicting Chapter 18 rather than smoothed over, which is
the behaviour that matters, since an instrument that only confirms is not an instrument.

The analysis is preregistered. \texttt{model/elicitation\_compare.py} was written before any form
came back, so what gets reported cannot be chosen after seeing the answers.

\section{What you cannot conclude from this}

Three respondents is not a validation study. Agreement would show that the table is reproducible
by people working from the same definitions, not that the definitions carve the world correctly.
The resource list is itself one person's, and question 2 on the form invites a respondent to say
it is wrong, which would be a more serious objection than any individual cell.

\end{document}
